\chapter{Ideals and the Splitting of Primes}
\begin{center}\itshape
In $\Z[\zeta_7]$ the prime $7$ becomes the sixth power of a single prime ideal, $29$ breaks into six different ones,\\
and $2$ into two. The rule behind this, and behind every number field, is the subject of this chapter.
\end{center}

On 1 March 1847, at a meeting of the Paris Academy of Sciences, Gabriel Lamé announced that he had proved Fermat's Last
Theorem. The idea was beautiful: factor $x^p+y^p$ over the $p$-th roots of unity, $x^p+y^p=\prod_k(x+\zeta^ky)$, and argue,
as one would with ordinary integers, that factors with no common divisor whose product is a $p$-th power must each be a
$p$-th power. Joseph Liouville stood up in the same session and asked the awkward question: who says factorisation into
primes is unique in $\Z[\zeta_p]$? Cauchy, sensing glory, announced that he too had been working on a proof. For weeks the
two filed sealed envelopes with the Academy.

The answer came from Breslau. Liouville received a letter from Ernst Kummer, read to the Academy in May, saying that
unique factorisation fails in $\Z[\zeta_{23}]$, that Kummer had published this three years earlier, and that he had
already repaired the damage. His repair was to invent ``ideal numbers'': phantom divisors that do not exist in the ring
but behave exactly like primes, so that every number factors uniquely into them. Twenty-four years later Dedekind made the
phantoms concrete. An ideal number is the set of all elements it ``divides'', and these sets are what we now call
ideals.

Tome~1, Chapter~11 told the quadratic part of this story, using a special trick: the conjugate ideal $\bar I$ and the
identity $I\bar I=(N(I))$. In higher degree there is no such conjugate, and we need a better reason why ideals factor.
Emmy Noether found it in 1927: three properties of $\mathcal O_K$ (it is Noetherian, integrally closed, and its non-zero
primes are maximal) are enough. Then we ask how a prime number $p$ breaks up in $\mathcal O_K$, and three characters
appear on stage. The numbers $e_i$ and $f_i$, with $\sum e_if_i=n$. The Kummer--Dedekind theorem, which reads the
factorisation off a polynomial modulo $p$. And, in Galois extensions, the star of the chapter: the \emph{Frobenius
element}, a single symmetry that ``is'' the map $x\mapsto x^p$. It explains why splitting in cyclotomic fields depends only
on $p\bmod m$, proves the theorem on Galois groups that Tome~2 left to Dedekind, and prepares Chebotarev's density theorem
(Chapter~8).

\begin{reminder}[from Tomes 1 and 2, and Chapter 1]
\emph{Ideals} (Tome~2, Propositions~6.5 and~6.8). In a commutative ring $R$, an ideal $I\ne R$ is prime if and only if
$R/I$ is an integral domain, and maximal if and only if $R/I$ is a field; a finite integral domain is a field (Tome~2,
Proposition~6.3). The product $IJ$ is the ideal generated by the products $ab$, $a\in I$, $b\in J$.
\emph{Chinese remainder theorem} (Tome~2, Theorem~6.9). If $I+J=R$, then $IJ=I\cap J$ and $R/IJ\cong R/I\times R/J$; by
induction the same holds for finitely many pairwise coprime ideals.
\emph{Finite fields} (Tome~2, Theorems~9.6 and~9.8). For each $q=p^f$ there is one field $\mathbb F_q$ up to isomorphism;
its automorphism group is cyclic of order $f$, generated by $\varphi:x\mapsto x^p$. If $\bar g\in\mathbb F_p[X]$ is
irreducible of degree $d$ with a root $\beta$, its roots are $\beta,\beta^p,\dots,\beta^{p^{d-1}}$, all distinct.
\emph{Galois theory} (Tome~2, Theorem~10.6, Proposition~10.9 and Exercise~10.8). The fundamental theorem; the cyclotomic
polynomial $\Phi_m$ is irreducible and $\operatorname{Gal}(\Q(\zeta_m)/\Q)\cong(\Z/m\Z)^*$, $\sigma_a(\zeta_m)=\zeta_m^a$.
\emph{Chapter~1.} Integrality is tested by finitely generated modules (Proposition~\ref{prop:intcrit}); $\mathcal O_K$
is a lattice with an integral basis (Theorem~\ref{thm:intbasis}); $N(\alpha)=\det\mu_\alpha$; the index formula
$\operatorname{disc}(\theta)=[\mathcal O_K:\Z[\theta]]^2d_K$ (Proposition~\ref{prop:index}).
\end{reminder}

\section{Three properties of $\mathcal O_K$}

Throughout, $K$ is a number field of degree $n$ and $\mathcal O=\mathcal O_K$. ``Ideal'' means non-zero ideal of
$\mathcal O$, and ``prime'' means non-zero prime ideal, unless we speak of prime numbers.

\begin{proposition}\label{prop:noeth}
\begin{enumerate}[label=(\alph*)]
\item Every ideal $I$ contains a positive integer, and $\mathcal O/I$ is finite.
\item Every ideal is a free abelian group of rank $n$, generated as an ideal by at most $n$ elements; and an ideal is
contained in only finitely many ideals.
\item $\mathcal O$ is \emph{integrally closed}: an element of $K$ that is a root of a monic polynomial with coefficients in
$\mathcal O$ lies in $\mathcal O$.
\item Every prime ideal is maximal, and its residue field $\mathcal O/P$ is finite.
\end{enumerate}
\end{proposition}
\begin{proof}
(a) Let $\alpha\in I$, $\alpha\ne0$. Then $N(\alpha)=\alpha\cdot\prod_{i\ge2}\sigma_i(\alpha)$, where $\sigma_1$ is the
identity. The product $\beta=\prod_{i\ge2}\sigma_i(\alpha)=N(\alpha)/\alpha$ lies in $K$, and it is an algebraic integer
as a product of conjugates of $\alpha$ (Proposition~\ref{prop:minpolyZ} and Corollary~\ref{cor:ring}). So $\beta\in\mathcal
O$ and $m=|N(\alpha)|=\pm\alpha\beta\in I$, a positive integer. Then $m\mathcal O\subseteq I$, and $\mathcal O/m\mathcal
O\cong(\Z/m\Z)^n$ by Theorem~\ref{thm:intbasis}, so $\mathcal O/I$, a quotient of it, is finite.

(b) $I$ is a subgroup of $\mathcal O\cong\Z^n$, hence free of rank at most $n$ (Tome~2, Lemma~7.5), and of rank $n$
since it contains $m\mathcal O\cong\Z^n$. A $\Z$-basis generates $I$ as an ideal. An ideal $J\supseteq I$ contains
$m\mathcal O$, so it is determined by its image in the finite ring $\mathcal O/m\mathcal O$.

(c) If $x\in K$ is a root of a monic polynomial with coefficients in $\mathcal O$, it is an algebraic integer by
Corollary~\ref{cor:transitive}, so $x\in\mathcal O$.

(d) If $P$ is prime, $\mathcal O/P$ is an integral domain, finite by (a), hence a field.
\end{proof}

Rings with these three properties (Noetherian, integrally closed, every non-zero prime maximal) are called
\emph{Dedekind domains}. Everything in the next section holds for them, and Tome~9 develops that abstract theory; we
stay with $\mathcal O_K$, where finiteness of quotients makes the arguments shorter.

\section{Unique factorisation of ideals}

The integers $\Z$ have an inverse $\frac1p$ of $p$ in $\Q$. Ideals need the analogue.

\begin{definition}
A \emph{fractional ideal} of $K$ is a set $J=\frac1d I$ where $I$ is an ideal of $\mathcal O$ and $d\in\Z$, $d\ge1$.
Products are defined as for ideals. For a prime $P$, put $P^{-1}=\{x\in K:xP\subseteq\mathcal O\}$.
\end{definition}

$P^{-1}$ is a fractional ideal: it contains $\mathcal O$, it is an $\mathcal O$-module, and if $m\in P$ is a positive
integer then $mP^{-1}\subseteq\mathcal O$. The whole point is that $P^{-1}$ is strictly larger than $\mathcal O$. To find an
element in it we need a lemma, whose proof is the first use of finiteness.

\begin{lemma}\label{lem:containsprod}
Every ideal $I$ contains a product $P_1\cdots P_r$ of primes ($r\ge0$; the empty product is $\mathcal O$, used only when
$I=\mathcal O$).
\end{lemma}
\begin{proof}
By strong induction on the finite number $[\mathcal O:I]$. If $I=\mathcal O$ there is nothing to do, and if $I$ is prime
we take $r=1$. Otherwise there are $a,b\in\mathcal O\setminus I$ with $ab\in I$. The ideals $I+a\mathcal O$ and
$I+b\mathcal O$ strictly contain $I$, so they have smaller index and contain products of primes. Their product
$(I+a\mathcal O)(I+b\mathcal O)$ is contained in $I+ab\mathcal O=I$, and it contains the product of the two products.
\end{proof}

\begin{lemma}[the missing element]\label{lem:missing}
For every prime $P$, $P^{-1}\ne\mathcal O$.
\end{lemma}
\begin{proof}
Take $a\in P$, $a\ne0$. By Lemma~\ref{lem:containsprod}, $a\mathcal O\supseteq P_1\cdots P_r$ with primes $P_i$; choose
such a product with $r$ as small as possible ($r\ge1$, since $a\mathcal O\subseteq P\ne\mathcal O$). Then
$P\supseteq P_1\cdots P_r$. A prime ideal containing a product of ideals contains one of them (if $x_i\in P_i\setminus P$
for all $i$, then $x_1\cdots x_r\in P$ while no $x_i$ is, which a prime ideal forbids). So $P\supseteq P_i$ for some $i$,
say $i=1$, and $P=P_1$ because $P_1$ is maximal. By minimality of $r$, $a\mathcal O\not\supseteq P_2\cdots P_r$: pick
$b\in P_2\cdots P_r$ with $b\notin a\mathcal O$. Then $bP\subseteq P_1P_2\cdots P_r\subseteq a\mathcal O$, so
$\frac baP\subseteq\mathcal O$, that is $\frac ba\in P^{-1}$; and $\frac ba\notin\mathcal O$.
\end{proof}

\begin{proposition}\label{prop:pinv}
For every prime $P$ and every ideal $I$: $IP^{-1}\ne I$, and $PP^{-1}=\mathcal O$.
\end{proposition}
\begin{proof}
Suppose $IP^{-1}=I$. Then every $x\in P^{-1}$ satisfies $xI\subseteq I$. As $I$ is a non-zero finitely generated additive
group (Proposition~\ref{prop:noeth}), $x$ is an algebraic integer by Proposition~\ref{prop:intcrit}(c), so $x\in\mathcal
O$. Hence $P^{-1}\subseteq\mathcal O$, contradicting Lemma~\ref{lem:missing}.

For the second claim, $P=P\mathcal O\subseteq PP^{-1}\subseteq\mathcal O$, and $PP^{-1}$ is an ideal. As $P$ is maximal,
$PP^{-1}$ is $P$ or $\mathcal O$, and it is not $P$ by the first claim with $I=P$.
\end{proof}

\begin{theorem}[Dedekind]\label{thm:ufi}
Every ideal $I\ne\mathcal O$ is a product of primes, $I=P_1\cdots P_r$, uniquely up to the order of the factors.
\end{theorem}
\begin{proof}
\emph{Existence}, by strong induction on $[\mathcal O:I]$. The ideals containing $I$ are finitely many; among those
different from $\mathcal O$ choose one, $P$, contained in no other. It is maximal, hence prime. Then
$I\subseteq IP^{-1}\subseteq PP^{-1}=\mathcal O$, so $J=IP^{-1}$ is an ideal, and $J\ne I$ by
Proposition~\ref{prop:pinv}. So $J\supsetneq I$ has smaller index, and $JP=IP^{-1}P=I$. Either $J=\mathcal O$ and $I=P$, or $J$ is a product of primes by induction, and so is $I=JP$.

\emph{Uniqueness.} Let $P_1\cdots P_r=Q_1\cdots Q_s$. Then $P_1$ contains $Q_1\cdots Q_s$, so it contains some $Q_j$ (as
in the proof of Lemma~\ref{lem:missing}), and $P_1=Q_j$ since $Q_j$ is maximal. Multiplying both sides by $P_1^{-1}$
and using $P_1P_1^{-1}=\mathcal O$, we cancel it, and conclude by induction on $r$. When $r=0$ the left side is
$\mathcal O$, while a product with $s\ge1$ factors is contained in $Q_1\ne\mathcal O$; so $s=0$.
\end{proof}

We write $v_P(I)$ for the exponent of $P$ in $I$, so that $I=\prod_PP^{v_P(I)}$, a finite product. For an element,
$v_P(\alpha)=v_P(\alpha\mathcal O)$.

\begin{corollary}\label{cor:divides}
\begin{enumerate}[label=(\alph*)]
\item The fractional ideals form a group under multiplication, free abelian on the primes: every fractional ideal is
uniquely $\prod_PP^{a_P}$ with $a_P\in\Z$, almost all zero.
\item (To contain is to divide.) For ideals, $I\subseteq J$ if and only if $I=JL$ for an ideal $L$, if and only if
$v_P(I)\ge v_P(J)$ for all $P$.
\item $I+J$ and $I\cap J$ are the ideals with exponents $\min(v_P(I),v_P(J))$ and $\max(v_P(I),v_P(J))$.
\end{enumerate}
\end{corollary}
\begin{proof}
(a) Put $P^{-a}=(P^{-1})^a$. Every fractional ideal is $\frac1dI$, and $\frac1d\mathcal O$ is the inverse of the ideal
$d\mathcal O$, a product of primes; so it is a product of primes and inverses of primes, which we may collect as
$\prod P^{a_P}$. Uniqueness: if $\prod P^{a_P}=\prod P^{b_P}$, multiply by a large product of primes to make all
exponents non-negative and apply Theorem~\ref{thm:ufi}. The inverse of $\prod P^{a_P}$ is $\prod P^{-a_P}$.

(b) If $I\subseteq J$, then $L=J^{-1}I\subseteq J^{-1}J=\mathcal O$ is an ideal, and $JL=I$. The converse is clear, and
the condition on exponents is a restatement by (a).

(c) $I+J$ is the smallest ideal containing $I$ and $J$, and $I\cap J$ the largest contained in both. By (b), these are
the ideals with the largest exponents $\le$ both, and the smallest $\ge$ both.
\end{proof}

So ideals behave like positive integers: ``contains'' is ``divides'', $+$ is $\gcd$ and $\cap$ is $\operatorname{lcm}$.
Two ideals are \emph{coprime} when $I+J=\mathcal O$, that is, when they have no prime factor in common. The
Chinese remainder theorem then applies to any distinct primes: $P^a$ and $Q^b$ are coprime, so
\begin{equation}\label{eq:crt}
\mathcal O/I\cong\prod_{P\mid I}\mathcal O/P^{v_P(I)}.
\end{equation}

\begin{definition}
The \emph{class group} $\operatorname{Cl}(K)$ is the quotient of the group of fractional ideals by the subgroup of
principal fractional ideals $\alpha\mathcal O$, $\alpha\in K^\times$. Its order $h_K$ is the \emph{class number}.
\end{definition}

For quadratic fields this agrees with Tome~1's definition (Tome~1, Section~11.6). Chapter~3 proves that $h_K$ is
always finite.

\section{The norm of an ideal}

\begin{definition}
The \emph{norm} of an ideal $I$ is $N(I)=[\mathcal O:I]$, finite by Proposition~\ref{prop:noeth}.
\end{definition}

\begin{proposition}\label{prop:normideal}
\begin{enumerate}[label=(\alph*)]
\item $N(IJ)=N(I)N(J)$ for all ideals $I,J$.
\item $N(\alpha\mathcal O)=|N(\alpha)|$ for $\alpha\in\mathcal O\setminus\{0\}$.
\item $N(I)\in I$. Each prime $P$ contains exactly one prime number $p$, and $N(P)=p^f$ with $f=[\mathcal O/P:\mathbb F_p]$.
\end{enumerate}
\end{proposition}
\begin{proof}
(a) First, $[I:IP]=N(P)$ for every ideal $I$ and prime $P$. The group $I/IP$ is an $\mathcal O$-module killed by $P$,
hence a vector space over the field $k=\mathcal O/P$, and its subspaces are the $L/IP$ with $L$ an ideal, $IP\subseteq
L\subseteq I$. For such $L$, $I^{-1}L$ is an ideal between $P$ and $\mathcal O$, hence equal to one of them, and $L$ is
$IP$ or $I$. So $I/IP$ has no subspace but $0$ and itself, and it is non-zero ($IP\ne I$ by cancellation): it has
dimension $1$, and $[I:IP]=|k|=N(P)$. Then $[\mathcal O:IP]=[\mathcal O:I][I:IP]=N(I)N(P)$, and by induction on the
number of prime factors of $J$, $N(IJ)=N(I)N(J)$. (This is Tome~2's proof of Theorem~7.3, with $P^{-1}$ in place of $\bar
P$.)

(b) $\alpha\mathcal O$ is the image of $\mathcal O$ under $\mu_\alpha$. In an integral basis, $\mu_\alpha$ has an integer
matrix with determinant $N(\alpha)\ne0$, and the index of the image is $|\det\mu_\alpha|$ (Tome~2, Corollary~7.8).

(c) The ring $\mathcal O/I$ has $N(I)$ elements, so $N(I)\cdot\bar1=0$ in it: $N(I)\in I$. A prime $P$ contains
$N(P)>1$, hence one of its prime factors $p$; if it contained two different prime numbers $p,q$, it would contain
$1=up+vq$. The field $\mathcal O/P$ contains $\Z/(P\cap\Z)=\mathbb F_p$, so it is a vector space of dimension
$f$ over it and has $p^f$ elements.
\end{proof}

By (c), a prime $P$ lies \emph{above} exactly one prime number $p$, and $P\mid p\mathcal O$. So to find all the primes of
$\mathcal O$ we factor each $p\mathcal O$. Write
\[
p\mathcal O=P_1^{e_1}\cdots P_g^{e_g},\qquad f_i=[\mathcal O/P_i:\mathbb F_p].
\]
The number $e_i$ is the \emph{ramification index} and $f_i$ the \emph{residue degree} of $P_i$.

\begin{theorem}\label{thm:sumef}
$\displaystyle\sum_{i=1}^ge_if_i=n$.
\end{theorem}
\begin{proof}
$N(p\mathcal O)=|N(p)|=p^n$ by Proposition~\ref{prop:normideal}(b), and $N(p\mathcal O)=\prod N(P_i)^{e_i}=p^{\sum
e_if_i}$ by (a) and (c).
\end{proof}

We say $p$ \emph{splits completely} if $g=n$ (all $e_i=f_i=1$), is \emph{inert} if $p\mathcal O$ is prime ($g=1$,
$f_1=n$), is \emph{totally ramified} if $p\mathcal O=P^n$, and \emph{ramifies} if some $e_i\ge2$. Figure~\ref{fig:spec}
shows the three behaviours in $\Z[i]$.

\begin{figure}[ht]
\centering
\fig{ch02-spec}
\caption{The primes of $\Z[i]$ drawn above the primes of $\Z$, in the style of Mumford's ``Red Book''. Above each $p$
lie the primes of $\Z[i]$ containing it: two for $p\equiv1\pmod 4$, one ``fat'' point of residue degree $2$ for
$p\equiv3\pmod4$, and the single point $(1+i)$ where the two branches meet, above the ramified prime $2$.}
\label{fig:spec}
\end{figure}

\section{The Kummer--Dedekind theorem}

How does one find the $P_i$? Tome~1 did it for quadratic fields with explicit generators $(p,a+\sqrt d)$. In general
the same recipe works whenever $\mathcal O$ looks like $\Z[\theta]$ near $p$.

\begin{theorem}[Kummer, Dedekind]\label{thm:kd}
Let $K=\Q(\theta)$ with $\theta\in\mathcal O$, $f=m_\theta$, and let $p$ be a prime number not dividing
$[\mathcal O:\Z[\theta]]$. Factor $f$ modulo $p$ as
\[
\bar f=\bar g_1^{e_1}\cdots\bar g_g^{e_g},
\]
with distinct monic irreducible $\bar g_i\in\mathbb F_p[X]$, and let $g_i\in\Z[X]$ be any monic lifts. Then the ideals
$P_i=p\mathcal O+g_i(\theta)\mathcal O$ are distinct primes, $f(P_i)=\deg\bar g_i$, and
\[
p\mathcal O=P_1^{e_1}\cdots P_g^{e_g}.
\]
\end{theorem}
\begin{proof}
\emph{Step 1: $\mathcal O/p\mathcal O\cong\mathbb F_p[X]/(\bar f)$.} Consider the ring map $\Z[\theta]\to\mathcal O/p
\mathcal O$. Its kernel is $\Z[\theta]\cap p\mathcal O$. If $p\beta\in\Z[\theta]$ with $\beta\in\mathcal O$, the class
of $\beta$ in the finite group $\mathcal O/\Z[\theta]$ has order dividing $p$, and this group has order prime to $p$; so
$\beta\in\Z[\theta]$. The kernel is therefore $p\Z[\theta]$, and the map induces an injection $\Z[\theta]/p\Z[\theta]\to
\mathcal O/p\mathcal O$. Both sides have $p^n$ elements ($\Z[\theta]$ and $\mathcal O$ are free of rank $n$), so it is an
isomorphism. Finally $\Z[\theta]\cong\Z[X]/(f)$, so $\Z[\theta]/p\Z[\theta]\cong\Z[X]/(p,f)\cong\mathbb F_p[X]/(\bar f)$.
The composite isomorphism sends the class of $h(X)$ to $h(\theta)\bmod p\mathcal O$.

\emph{Step 2: the primes above $p$.} The primes of $\mathcal O$ containing $p$ correspond to the maximal ideals of
$\mathcal O/p\mathcal O$, hence to those of $\mathbb F_p[X]/(\bar f)$, which are the $(\bar g_i)/(\bar f)$ (a maximal
ideal of $\mathbb F_p[X]$ containing $\bar f$ is generated by an irreducible factor of $\bar f$). The prime of
$\mathcal O$ corresponding to $(\bar g_i)$ is the inverse image of the ideal generated by $g_i(\theta)$, which is
$P_i=p\mathcal O+g_i(\theta)\mathcal O$. So the $P_i$ are distinct primes, they are all the primes above $p$, and
$\mathcal O/P_i\cong\mathbb F_p[X]/(\bar g_i)$ is a field with $p^{\deg g_i}$ elements.

\emph{Step 3: the exponents.} The ideal $P_1^{e_1}\cdots P_g^{e_g}$ is generated by $p$ and elements that are products of
factors $p$ and $g_i(\theta)$; modulo $p\mathcal O$ all of them are multiples of $\prod g_i(\theta)^{e_i}\equiv
f(\theta)=0$. So $\prod P_i^{e_i}\subseteq p\mathcal O$, hence $p\mathcal O\mid\prod P_i^{e_i}$
(Corollary~\ref{cor:divides}), and $p\mathcal O=\prod P_i^{a_i}$ with $a_i\le e_i$ (only the $P_i$ divide $p\mathcal O$).
Taking norms, $p^n=p^{\sum a_i\deg g_i}$, while $\sum e_i\deg g_i=\deg f=n$. So $\sum(e_i-a_i)\deg g_i=0$ and $a_i=e_i$.
\end{proof}

The theorem applies to all but finitely many primes, since the index has finitely many prime factors. It reduces the
factorisation of $p$ to the factorisation of a polynomial over a finite field, a fast computation using, for example, Cantor and Zassenhaus' algorithm.

\begin{example}[{$\Q(\sqrt[3]2)$}]\label{exa:cubesplit}
Here $\mathcal O=\Z[\sqrt[3]2]$ (Example~\ref{exa:cuberoot2}), so the theorem applies to every $p$. Modulo $2$,
$X^3-2\equiv X^3$, so $2\mathcal O=(\sqrt[3]2)^3$: totally ramified. Modulo $3$, $X^3-2\equiv(X+1)^3$, so $3\mathcal
O=(3,\sqrt[3]2+1)^3$, and since $N(\sqrt[3]2+1)=3$ this prime is the principal ideal $(\sqrt[3]2+1)$. Modulo $5$, $X^3-2
\equiv(X-3)(X^2+3X+4)$, with the quadratic factor irreducible (its discriminant $-7\equiv3$ is not a square mod $5$):
$5\mathcal O=P_1P_2$ with $f_1=1$, $f_2=2$. Modulo $7$, $2$ is not a cube ($\pm1$ and $0$ are the only cubes), so $7$ is
inert. Modulo $31$, $2\equiv4^3$ and $X^3-2$ has the three roots $4,7,20$: $31$ splits completely.
Exercise~\ref{ex:cubic} finds the rule.
\end{example}

\begin{example}[{the prime cyclotomic field}]\label{exa:zetap}
Let $\ell$ be an odd prime, $\zeta=\zeta_\ell$, $K=\Q(\zeta)$, of degree $\ell-1$. Differentiating $(X-1)\Phi_\ell(X)=
X^\ell-1$ at $\zeta^a$ gives $(\zeta^a-1)\Phi_\ell'(\zeta^a)=\ell\zeta^{-a}$, and taking the product over $a=1,\dots,\ell-1$,
with $\prod_a(1-\zeta^a)=\Phi_\ell(1)=\ell$, we find $N(\Phi_\ell'(\zeta))=\pm\ell^{\ell-1}/\ell=\pm\ell^{\ell-2}$. So
$\operatorname{disc}(\Phi_\ell)=\pm\ell^{\ell-2}$ by Proposition~\ref{prop:disc}(c), and only $\ell$ can divide the index
of $\Z[\zeta]$. But $\Phi_\ell(X+1)=\frac{(X+1)^\ell-1}X$ is $\ell$-Eisenstein, so $\ell$ does not divide it
(Proposition~\ref{prop:eisenstein}). Hence $\mathcal O_K=\Z[\zeta_\ell]$ and $d_K=\pm\ell^{\ell-2}$. (Chapter~5 extends this
to all $\Q(\zeta_m)$.)

Now Theorem~\ref{thm:kd} applies to every $p$. For $p=\ell$, $\Phi_\ell\equiv\frac{X^\ell-1}{X-1}=(X-1)^{\ell-1}$, so
$\ell\mathcal O=(\ell,\zeta-1)^{\ell-1}=(1-\zeta)^{\ell-1}$: totally ramified, as $N(1-\zeta)=\Phi_\ell(1)=\ell$ shows that
$(1-\zeta)$ is already the prime. For $\ell=7$ and $p=29\equiv1\pmod7$, $\mathbb F_{29}^*$ is cyclic of order $28$ and
contains the six primitive $7$-th roots of unity $16,24,7,25,23,20$ (the powers of $16=2^{4}$, as $2$ generates
$\mathbb F_{29}^*$): $\Phi_7$ splits into six linear factors, and so does $29$. For $p=13$, the order of $13$ modulo $7$ is
$2$; Theorem~\ref{thm:cyclosplit} below shows that this order is the residue degree.
\end{example}

\begin{think}
Factor $\Phi_7$ modulo $2$ by hand (try $X^3+X+1$). How many primes lie above $2$ in $\Z[\zeta_7]$, and of which degree?
Guess the rule relating the factorisation of $p$ to the order of $p$ modulo $7$.
\end{think}

\section{Ramification and the discriminant}

\begin{theorem}\label{thm:ramdisc}
A prime number $p$ ramifies in $K$ if and only if $p\mid d_K$. In particular only finitely many primes ramify.
\end{theorem}
\begin{proof}
Let $A=\mathcal O/p\mathcal O$, an algebra of dimension $n$ over $\mathbb F_p$ with basis the classes $\bar\omega_j$ of an
integral basis. For $x\in\mathcal O$, the matrix of $\mu_x$ in the basis $(\omega_j)$ has integer entries, and its
reduction modulo $p$ is the matrix of multiplication by $\bar x$ on $A$. Taking traces, $\operatorname{Tr}_{K/\Q}(x)\bmod
p=\operatorname{Tr}_A(\bar x)$, so by Proposition~\ref{prop:disc}(a)
\[
d_K\bmod p=\det\bigl(\operatorname{Tr}_A(\bar\omega_i\bar\omega_j)\bigr).
\]
This is the determinant of the \emph{trace form} $(u,v)\mapsto\operatorname{Tr}_A(uv)$ of $A$ in a basis. So $p\mid
d_K$ if and only if the trace form of $A$ is degenerate, a property that does not depend on the basis.

By \eqref{eq:crt}, $A\cong B_1\times\dots\times B_g$ with $B_i=\mathcal O/P_i^{e_i}$. The trace of multiplication by
$(b_1,\dots,b_g)$ on a product is the sum of the traces on the factors, so in a basis adapted to the product the matrix
of the trace form is block diagonal, and $A$'s form is non-degenerate if and only if each $B_i$'s is.

If $e_i\ge2$, pick $x\in P_i^{e_i-1}\setminus P_i^{e_i}$. Its class $\bar x\in B_i$ is non-zero, and for every $y$,
$(\bar x\bar y)^{2}\in P_i^{2e_i-2}/P_i^{e_i}=0$ since $2e_i-2\ge e_i$. A nilpotent element has trace $0$ (a nilpotent
matrix has only the eigenvalue $0$). So $\operatorname{Tr}(\bar x\bar y)=0$ for all $\bar y$: the form is degenerate.

If $e_i=1$, $B_i=\mathcal O/P_i$ is the field $\mathbb F_{q}$, $q=p^{f_i}$, and multiplication by $u$ has trace
$u+u^p+\dots+u^{p^{f_i-1}}$. Indeed, let $\beta$ generate $\mathbb F_q^*$, so $\mathbb F_q=\mathbb F_p(\beta)$ with basis
$1,\beta,\dots,\beta^{f_i-1}$, and $u=h(\beta)$ with $h\in\mathbb F_p[X]$. The minimal polynomial of $\beta$ has the $f_i$
distinct roots $\beta^{p^k}$, $0\le k<f_i$ (Tome~2, Theorem~9.8), so $\mu_\beta$ is diagonalisable over $\mathbb F_q$ with
these eigenvalues, and $\mu_u=h(\mu_\beta)$ has eigenvalues $h(\beta^{p^k})=u^{p^k}$. This polynomial has degree $p^{f_i-1}<q$, so it does not vanish on all of
$\mathbb F_q$: some $u_0$ has non-zero trace. Then for $x\ne0$, $y=x^{-1}u_0$ gives $\operatorname{Tr}(xy)\ne0$: the form is
non-degenerate.

So $A$'s trace form is non-degenerate if and only if every $e_i=1$, which is the theorem.
\end{proof}

Together with Theorem~\ref{thm:kd} this gives a quick test. For $\Q(\sqrt d)$ with $d$ squarefree, $d_K=d$ or $4d$, and
the ramified primes are those dividing $d$, together with $2$ when $d\equiv2,3\pmod4$, as in Tome~1. For
$\Q(\sqrt[3]{10})$, $d_K=-300$: exactly $2$, $3$ and $5$ ramify. Chapter~3 adds that $|d_K|>1$ for $K\ne\Q$: some prime
always ramifies.

\section{Galois extensions: decomposition, inertia, Frobenius}

Suppose now that $K/\Q$ is Galois with group $G$. Every $\sigma\in G$ maps $\mathcal O$ onto itself (conjugates of
algebraic integers are algebraic integers), hence maps primes to primes, and if $P$ lies above $p$ so does $\sigma(P)$,
because $\sigma(P)\cap\Z=\sigma(P\cap\Z)=p\Z$.

\begin{theorem}\label{thm:transitive}
$G$ acts transitively on the primes above $p$. Consequently all of them have the same $e$ and the same $f$, and
$efg=n$.
\end{theorem}
\begin{proof}
Let $P,Q$ be primes above $p$ and suppose $Q\ne\sigma(P)$ for all $\sigma\in G$. The primes $Q$ and $\sigma(P)$, $\sigma\in
G$, are pairwise coprime or equal, so by the Chinese remainder theorem there is $x\in\mathcal O$ with $x\equiv0\pmod Q$
and $x\equiv1\pmod{\sigma(P)}$ for every $\sigma$. Then $N(x)=\prod_{\tau\in G}\tau(x)$ lies in $Q\cap\Z=p\Z\subseteq P$.
As $P$ is prime, some $\tau(x)\in P$, that is $x\in\tau^{-1}(P)$, contradicting $x\equiv1\pmod{\tau^{-1}(P)}$.

If $Q=\sigma(P)$, applying $\sigma$ to $p\mathcal O=\prod P_i^{e_i}$ (which $\sigma$ fixes) permutes the factors and
shows $e(Q)=e(P)$; and $\sigma$ induces an isomorphism $\mathcal O/P\cong\mathcal O/Q$, so $f(Q)=f(P)$. Then
Theorem~\ref{thm:sumef} reads $efg=n$.
\end{proof}

\begin{definition}
Let $P$ be a prime above $p$, and $k_P=\mathcal O/P$. The \emph{decomposition group} is $D_P=\{\sigma\in G:\sigma(P)=P\}$.
Each $\sigma\in D_P$ induces an automorphism $\bar\sigma$ of $k_P$, $\bar\sigma(\bar x)=\overline{\sigma(x)}$. The
\emph{inertia group} $I_P$ is the kernel of $\sigma\mapsto\bar\sigma$, $D_P\to\operatorname{Gal}(k_P/\mathbb F_p)$.
\end{definition}

By transitivity, $D_P$ is the stabiliser of $P$ in an action with $g$ elements in the orbit, so $|D_P|=n/g=ef$. Also
$D_{\sigma(P)}=\sigma D_P\sigma^{-1}$ and $I_{\sigma(P)}=\sigma I_P\sigma^{-1}$.

\begin{theorem}\label{thm:decomp}
The map $D_P\to\operatorname{Gal}(k_P/\mathbb F_p)$ is onto. Hence $|I_P|=e$, and if $p$ is unramified, $D_P$ is cyclic of
order $f$, generated by a unique $\operatorname{Frob}_P\in D_P$ such that
\[
\operatorname{Frob}_P(x)\equiv x^p\pmod P\qquad\text{for all }x\in\mathcal O.
\]
\end{theorem}
\begin{proof}
The multiplicative group $k_P^*$ is cyclic (Tome~2, Theorem~9.4); let $\bar a$ be a generator, so $k_P=\mathbb F_p(\bar
a)$. By the Chinese remainder theorem choose $a\in\mathcal O$ with $a\equiv\bar a\pmod P$ and $a\equiv0\pmod
{\sigma^{-1}(P)}$ for every $\sigma\notin D_P$ (for those $\sigma$, $\sigma^{-1}(P)\ne P$). Consider
\[
h(X)=\prod_{\sigma\in G}\bigl(X-\sigma(a)\bigr),
\]
whose coefficients are fixed by $G$ and lie in $\mathcal O$, hence in $\Z$. Its reduction $\bar h\in\mathbb F_p[X]$ is
$\prod_\sigma(X-\overline{\sigma(a)})$ in $k_P[X]$, and it has $\bar a$ as a root. Let $\tau\in
\operatorname{Gal}(k_P/\mathbb F_p)$. Since $\bar h$ has coefficients in $\mathbb F_p$, $\tau(\bar a)$ is also a root, so
$\tau(\bar a)=\overline{\sigma(a)}$ for some $\sigma\in G$. If $\sigma\notin D_P$, then $a\in\sigma^{-1}(P)$, so
$\sigma(a)\in P$ and $\overline{\sigma(a)}=0\ne\tau(\bar a)$. So $\sigma\in D_P$, and $\bar\sigma(\bar a)=\tau(\bar a)$.
As $\bar a$ generates $k_P$, $\bar\sigma=\tau$.

So $D_P/I_P\cong\operatorname{Gal}(k_P/\mathbb F_p)$, of order $f$, and $|I_P|=ef/f=e$. If $e=1$, $D_P\cong
\operatorname{Gal}(k_P/\mathbb F_p)$, which is cyclic, generated by $x\mapsto x^p$ (Tome~2, Theorem~9.8); its preimage is
$\operatorname{Frob}_P$.
\end{proof}

The Frobenius elements above $p$ are conjugate: $\operatorname{Frob}_{\sigma(P)}=\sigma\operatorname{Frob}_P\sigma^{-1}$,
because the right side fixes $\sigma(P)$ and reduces to $x\mapsto x^p$ modulo $\sigma(P)$. So an unramified $p$ has a
well-defined conjugacy class $\operatorname{Frob}_p\subseteq G$, the \emph{Frobenius class}, and if $G$ is abelian, a
well-defined element. Its order is $f$, so $p$ splits completely if and only if $\operatorname{Frob}_p=1$.

\begin{think}
Why can a prime never be inert in a Galois extension whose group is not cyclic? What does this say about $X^4+1$
modulo $p$, whose splitting field $\Q(\zeta_8)$ has group $(\Z/8\Z)^*\cong(\Z/2\Z)^2$?
\end{think}

\begin{theorem}[splitting in cyclotomic fields]\label{thm:cyclosplit}
Let $K=\Q(\zeta_m)$ and $p\nmid m$. Then $p$ is unramified in $K$, $\operatorname{Frob}_p=\sigma_p:\zeta_m\mapsto
\zeta_m^p$, and $p$ splits into $g=\varphi(m)/f$ primes of residue degree $f$, the order of $p$ in $(\Z/m\Z)^*$.
\end{theorem}
\begin{proof}
\emph{Unramified.} Write $X^m-1=\Phi_m(X)H(X)$ in $\Z[X]$. Differentiating at $\zeta=\zeta_m$ gives
$\Phi_m'(\zeta)H(\zeta)=m\zeta^{m-1}$, and taking norms, $N(\Phi_m'(\zeta))\,N(H(\zeta))=\pm m^{\varphi(m)}$, with
$N(H(\zeta))$ a non-zero integer (the roots of $X^m-1$ are distinct, so $H(\zeta)\ne0$). So
$\operatorname{disc}(\Phi_m)=\pm N(\Phi_m'(\zeta))$ divides $m^{\varphi(m)}$, and $p\nmid\operatorname{disc}(\Phi_m)=
[\mathcal O:\Z[\zeta]]^2d_K$. So $p\nmid d_K$ (Theorem~\ref{thm:ramdisc}).

\emph{Frobenius.} In a commutative ring, $(u_1+\dots+u_k)^p\equiv u_1^p+\dots+u_k^p\pmod p$, since the multinomial
coefficients other than those of the $u_i^p$ are divisible by $p$. So for $y=\sum_ia_i\zeta^i$ with $a_i\in\Z$, Fermat's
little theorem gives $y^p\equiv\sum_ia_i^p\zeta^{ip}\equiv\sum_ia_i\zeta^{ip}=\sigma_p(y)\pmod{p\mathcal O}$. We do not
know yet that $\mathcal O=\Z[\zeta]$, but $p\nmid[\mathcal O:\Z[\zeta]]$, so every $x\in\mathcal O$ is $y+pz$ with
$y\in\Z[\zeta]$, $z\in\mathcal O$ (Step~1 of Theorem~\ref{thm:kd}). Then $x^p\equiv y^p\pmod{p\mathcal O}$ and
$\sigma_p(x)=\sigma_p(y)+p\sigma_p(z)\equiv\sigma_p(y)$, as $\sigma_p(\mathcal O)=\mathcal O$. So $\sigma_p(x)\equiv x^p
\pmod{p\mathcal O}$ for all $x\in\mathcal O$. If $x\in P$,
then $\sigma_p(x)\equiv x^p\equiv0\pmod P$, so $\sigma_p(P)\subseteq P$, hence $\sigma_p\in D_P$, and it reduces to
$x\mapsto x^p$: $\sigma_p=\operatorname{Frob}_P$ for every $P$ above $p$.

\emph{Splitting.} The order of $\sigma_p$ in $G\cong(\Z/m\Z)^*$ is the order of $p$ modulo $m$, and it is $f$; then
$g=\varphi(m)/f$ by Theorem~\ref{thm:transitive}.
\end{proof}

So the splitting of $p$ in $\Q(\zeta_m)$ depends only on $p\bmod m$: this is the simplest instance of a \emph{reciprocity
law}, and class field theory (Chapter~10) says that every abelian extension of $\Q$ behaves this way. For the opening
example, $13\equiv-1\pmod 7$ has order $2$, so $13$ splits into $6/2=3$ primes of degree $2$ in $\Z[\zeta_7]$; a prime
splits into exactly two primes, of degree $3$, when its order modulo $7$ is $3$, for instance $p=2$ or $p=11$.

\section{Dedekind's theorem on Galois groups}

Here is the theorem that Tome~2 (Chapter~10, the remark on Frobenius and reduction modulo $p$) used to guess Galois groups from factorisations modulo $p$.

\begin{theorem}[Dedekind]\label{thm:dedekindcycle}
Let $f\in\Z[X]$ be monic and irreducible of degree $n$, with roots $\theta_1,\dots,\theta_n$, splitting field $L$, and
Galois group $G\subseteq S_n$ acting on the roots. Let $p$ be a prime with $p\nmid\operatorname{disc}(f)$, and suppose that
$\bar f$ factors in $\mathbb F_p[X]$ into irreducible factors of degrees $d_1,\dots,d_r$. Then $G$ contains a permutation
of cycle type $(d_1,\dots,d_r)$; it is the Frobenius of any prime of $L$ above $p$.
\end{theorem}
\begin{proof}
Let $P$ be a prime of $\mathcal O_L$ above $p$ and $\sigma\in D_P$ with $\bar\sigma(\bar x)=\bar x^p$ on $k_P$, which
exists by Theorem~\ref{thm:decomp}. The roots $\theta_i$ are in $\mathcal O_L$. Since $\operatorname{disc}(f)=\prod_{i<j}
(\theta_i-\theta_j)^2$ is an integer not divisible by $p$, it is not in $P$, so the reductions $\bar\theta_i\in k_P$ are
pairwise distinct; and $\bar f=\prod_i(X-\bar\theta_i)$ in $k_P[X]$. Now $\sigma(\theta_i)=\theta_j$ for some $j$, and
reducing, $\bar\theta_j=\bar\theta_i^{\,p}$. So, under the bijection $\theta_i\mapsto\bar\theta_i$, the permutation
$\sigma$ of the roots of $f$ becomes the permutation $\beta\mapsto\beta^p$ of the roots of $\bar f$. The roots of an
irreducible factor $\bar g$ of degree $d$ form one orbit $\beta,\beta^p,\dots,\beta^{p^{d-1}}$ of this map, of size $d$.
So the cycle type of $\sigma$ is $(d_1,\dots,d_r)$.
\end{proof}

\begin{example}
For $f=X^5-X-1$, $\operatorname{disc}(f)=2869=19\cdot151$. Modulo $5$, $f$ is irreducible (Exercise~\ref{ex:s5}), and
modulo $2$, $f\equiv(X^2+X+1)(X^3+X^2+1)$. So $G\subseteq S_5$ contains a $5$-cycle and a permutation $\sigma$ of type
$(2,3)$, whose cube is a transposition. A $5$-cycle and a transposition generate $S_5$ (as in the proof of Tome~2,
Theorem~10.14). So $\operatorname{Gal}(f)=S_5$, and $f$ is not solvable by radicals.
\end{example}

The same argument describes the primes of any number field through its Galois closure.

\begin{proposition}\label{prop:frobembed}
Let $K$ be a number field of degree $n$, $L$ its Galois closure, $P$ a prime of $L$ above $p$, and $\sigma\in D_P$ with
$\sigma(x)\equiv x^p\pmod P$ for $x\in\mathcal O_L$. Suppose $p$ is unramified in $K$, with $p\mathcal O_K=\mathfrak p_1
\cdots\mathfrak p_g$ of residue degrees $f_1,\dots,f_g$. Then $\tau\mapsto\sigma\circ\tau$ permutes the $n$ embeddings
$\tau:K\to L$ with cycle type $(f_1,\dots,f_g)$.
\end{proposition}
\begin{proof}
Let $A=\mathcal O_K/p\mathcal O_K\cong\prod_i\mathcal O_K/\mathfrak p_i=\prod_i\mathbb F_{p^{f_i}}$, and $k=\mathcal O_L/P$.
Each embedding $\tau$ maps $\mathcal O_K$ into $\mathcal O_L$, and reduction modulo $P$ gives a ring map $\bar\tau:A\to
k$. These $n$ maps are distinct: with an integral basis $(\omega_j)$, $\det(\tau(\omega_j))^2=d_K$ is not divisible by
$p$ (Theorem~\ref{thm:ramdisc}), so $\det(\bar\tau(\bar\omega_j))\ne0$ in $k$ and the rows are distinct. On the other
hand a ring map $A\to k$ is zero on all factors but one (it sends the idempotents of the product to idempotents of
the field $k$, that is to $0$ or $1$, with sum $1$ and pairwise products $0$), and on $\mathbb F_{p^{f_i}}$ it is one of at most $f_i$ field
embeddings (it is determined by the image of a generator $\beta$, a root of its minimal polynomial of degree $f_i$). So
there are at most $\sum f_i=n$ ring maps $A\to k$, and the $\bar\tau$ are all of them. Now $\overline{\sigma\circ\tau}=
\varphi\circ\bar\tau$ with $\varphi(y)=y^p$. The maps that are non-zero on the factor $\mathbb F_{p^{f_i}}$ are the
$\varphi^k\circ\psi$, $0\le k<f_i$, for one of them $\psi$, and $\varphi^{f_i}\circ\psi=\psi$; so they form one cycle of
length $f_i$.
\end{proof}

For $K=\Q(\theta)$, $p\nmid\operatorname{disc}(m_\theta)$, the embeddings correspond to the roots of $m_\theta$, and
Theorem~\ref{thm:kd} identifies the $f_i$ with the degrees of the factors of $m_\theta$ modulo $p$: this recovers
Theorem~\ref{thm:dedekindcycle}.

Dedekind's theorem says that each cycle type appearing modulo some $p$ appears in $G$. Chebotarev's theorem, proved in
Chapter~8, says much more: every conjugacy class $C$ of $G$ is the Frobenius class of a set of primes of density
$|C|/|G|$. Figure~\ref{fig:cycletypes} shows it at work for $X^5-X-1$.

\begin{figure}[ht]
\centering
\fig{ch02-cycletypes}
\caption{For the $17\,982$ primes $p<2\cdot10^5$ not dividing $2869$, the degrees of the irreducible factors of
$X^5-X-1$ modulo $p$ (dark bars), against the proportion of permutations of $S_5$ with that cycle type (light bars):
$\frac1{120}$, $\frac{10}{120}$, $\frac{20}{120}$, $\frac{15}{120}$, $\frac{30}{120}$, $\frac{20}{120}$,
$\frac{24}{120}$. The largest deviation is $0.003$. This is Chebotarev's density theorem (Chapter~8).}
\label{fig:cycletypes}
\end{figure}

\begin{think}
A polynomial of degree $5$ with Galois group $S_5$ has a root modulo $p$ for what proportion of primes? Count the
permutations of $S_5$ with at least one fixed point. Then do the same for a polynomial with group $A_5$.
\end{think}

\section{Exercises}

\begin{exercise}[$\star$]\label{ex:sqrtm10}
\emph{Two factorisations in $\Z[\sqrt{-10}]$.} Show that $14=2\cdot7=(2+\sqrt{-10})(2-\sqrt{-10})$ are two
factorisations into irreducible elements. Factor $2$, $7$, $2\pm\sqrt{-10}$ into prime ideals with
Theorem~\ref{thm:kd}, and see the two factorisations as two groupings of the same four prime ideals. Show that
$P=(2,\sqrt{-10})$ is not principal but $P^2$ is.
\end{exercise}

\begin{exercise}[$\star$]\label{ex:cubic}
Let $K=\Q(\sqrt[3]2)$ and $p\ge5$. Show that if $p\equiv2\pmod 3$, then $p=P_1P_2$ with $f_1=1$, $f_2=2$. If
$p\equiv1\pmod3$, show that $p$ splits completely or is inert, according as $2$ is or is not a cube modulo $p$.
\end{exercise}

\begin{exercise}[$\star$]\label{ex:zeta8}
Describe the splitting of every odd prime $p$ in $\Q(\zeta_8)=\Q(i,\sqrt2)$ according to $p\bmod8$. Deduce again that
$X^4+1$ is reducible modulo every prime, and find which of $\Q(i)$, $\Q(\sqrt2)$, $\Q(\sqrt{-2})$ the prime $p$ splits
in.
\end{exercise}

\begin{exercise}[$\star\star$]\label{ex:s5}
Show that $X^5-X-1$ is irreducible modulo $5$. \emph{(If $\beta$ is a root in some extension, so is $\beta+1$.)}
Check the factorisation modulo $2$ used in the text, and that $2869=19\cdot151$.
\end{exercise}

\begin{exercise}[$\star\star$]\label{ex:twogen}
\emph{Two generators suffice.} Let $I$ be an ideal and $a\in I$, $a\ne0$. Show that there is $b\in I$ with $I=a\mathcal
O+b\mathcal O$. \emph{(Find $b$ with $v_P(b)=v_P(I)$ for each $P\mid a\mathcal O$.)}
\end{exercise}

\begin{exercise}[$\star\star$]\label{ex:totram}
Let $q=\ell^k$ be a prime power and $\zeta=\zeta_q$. Show that $\ell\,\mathcal O_{\Q(\zeta)}=(1-\zeta)^{\varphi(q)}$ and
that $(1-\zeta)$ is a prime of residue degree $1$. \emph{(Compute $\Phi_q(1)$, and compare $1-\zeta^a$ with $1-\zeta$.)}
\end{exercise}

\begin{exercise}[$\star\star$]\label{ex:dedekindfield}
\emph{Dedekind's field again.} Let $K=\Q(\theta)$, $\theta^3-\theta^2-2\theta-8=0$, with integral basis $1,\theta,\beta$,
$\beta=\frac{\theta^2+\theta}2$ (Exercise~\ref{ex:dedekind}). Show that $\theta^2\equiv\theta$, $\beta^2\equiv\beta$ and
$\theta\beta\equiv0$ modulo $2\mathcal O$, deduce that $\mathcal O/2\mathcal O\cong\mathbb F_2^3$ and that $2$ splits
completely, and conclude that $\mathcal O\ne\Z[\gamma]$ for every $\gamma\in\mathcal O$.
\end{exercise}

\begin{exercise}[$\star\star$]\label{ex:closure}
Let $K$ be a number field and $L$ its Galois closure. Show that a prime $p$ splits completely in $K$ if and only if it
splits completely in $L$.
\end{exercise}

\begin{exercise}[$\star\star$]\label{ex:frobsub}
Let $L/\Q$ be Galois, $K\subseteq L$ a subfield with $K/\Q$ Galois, $P$ a prime of $L$ above an unramified $p$, and
$\mathfrak p=P\cap K$. Show that $\operatorname{Frob}_P$ restricts to $\operatorname{Frob}_{\mathfrak p}$ on $K$.
\end{exercise}

\begin{exercise}[$\star\star\star$]\label{ex:qrsplit}
\emph{Quadratic reciprocity from splitting.} Let $q$ be an odd prime, $q^*=(-1)^{(q-1)/2}q$, and $p\ne q$ an odd prime.
\begin{enumerate}[label=(\alph*)]
\item Show that $\Q(\zeta_q)$ has exactly one quadratic subfield, and that it is $\Q(\sqrt{q^*})$. \emph{(Tome~1,
Exercise~6.7: the Gauss sum.)}
\item Show that $p$ splits in $\Q(\sqrt{q^*})$ if and only if $\left(\frac{q^*}p\right)=1$.
\item Show that $p$ splits in $\Q(\sqrt{q^*})$ if and only if $p$ is a square modulo $q$.
\item Deduce the law of quadratic reciprocity.
\end{enumerate}
\end{exercise}

\begin{exercise}[$\star\star\star$]\label{ex:stickelbergerp}
\emph{Stickelberger's congruence for $p$.} Let $K$ be a number field of degree $n$ and $p$ an odd prime with $p\nmid d_K$,
with $g$ primes above it. Show that $\left(\frac{d_K}p\right)=(-1)^{n-g}$. Check it on $\Q(\sqrt[3]2)$ for $p=5,7,31$.
\end{exercise}

\begin{beauty}
Kummer invented ideal numbers in 1847 to repair unique factorisation in $\Z[\zeta_p]$, and Dedekind turned them into sets
in 1871; Noether's 1927 axioms (the three properties of Proposition~\ref{prop:noeth}) explained why the theory works. But the
deepest discovery of the chapter is the Frobenius element. It converts a question about one prime, how does $p$
factor, into a question about a group, which conjugacy class is $\operatorname{Frob}_p$; and for cyclotomic fields the
answer is just $p\bmod m$. Which number fields have such a ``modulus'' for their splitting? Exactly the abelian ones
(Chapter~10), and finding the analogue for non-abelian groups is the Langlands programme. Already for $\Q(\sqrt[3]2)$,
whose closure has group $S_3$, the rule ``$31$ splits because $2$ is a cube mod $31$'' has no congruence description;
Gauss's answer, $p=x^2+27y^2$, is proved in Chapter~10 and is the first non-abelian reciprocity law.
\end{beauty}

\section{Hints}

Try each exercise for a while before reading its hint, and read the hint before the solution.

\begin{hint}{ex:sqrtm10}
$\Z[\sqrt{-10}]$ is $\mathcal O$ since $-10\equiv2\pmod4$. Factor $X^2+10$ modulo $2$ and $7$. A generator of $P$ would
have norm $2=a^2+10b^2$.
\end{hint}
\begin{hint}{ex:cubic}
For $p\equiv2\pmod3$, cubing is a bijection of $\mathbb F_p^*$. For $p\equiv1\pmod3$, $\mathbb F_p$ contains the cube roots
of unity.
\end{hint}
\begin{hint}{ex:zeta8}
Use Theorem~\ref{thm:cyclosplit} with $m=8$: the order of $p$ modulo $8$ is $1$ or $2$. The three quadratic subfields
correspond to the three subgroups of order $2$ of $(\Z/8\Z)^*$.
\end{hint}
\begin{hint}{ex:s5}
$(\beta+1)^5-(\beta+1)-1=\beta^5-\beta-1$ in characteristic $5$. So the roots are $\beta,\beta+1,\dots,\beta+4$, and
the Frobenius permutes them as a translation.
\end{hint}
\begin{hint}{ex:twogen}
Use the Chinese remainder theorem with the ideals $P^{v_P(I)+1}$, $P\mid a\mathcal O$.
\end{hint}
\begin{hint}{ex:totram}
$\Phi_q(X)=\Phi_\ell(X^{q/\ell})$, so $\Phi_q(1)=\ell$. The quotient $\frac{1-\zeta^a}{1-\zeta}$ is in $\Z[\zeta]$, and so is
its inverse when $\gcd(a,\ell)=1$.
\end{hint}
\begin{hint}{ex:dedekindfield}
Use the multiplication table: $\theta^2=2\beta-\theta$, $\beta^2=6+2\theta+3\beta$, $\theta\beta=4+2\beta$. In a
finite ring where $x^2=x$ for all $x$, the idempotents split the ring into a product. If $\mathcal O=\Z[\gamma]$, then
$\mathcal O/2\mathcal O\cong\mathbb F_2[X]/(\bar m_\gamma)$.
\end{hint}
\begin{hint}{ex:closure}
Proposition~\ref{prop:frobembed}: if $p$ splits completely in $K$, a Frobenius $\sigma$ fixes every embedding of $K$, so
it is trivial on every conjugate field $\tau(K)$, which generate $L$. For ramification, look for nilpotent elements of
$\mathcal O/p\mathcal O$.
\end{hint}
\begin{hint}{ex:frobsub}
Restriction maps $D_P$ into $D_{\mathfrak p}$, and the reduction of $x\in\mathcal O_K$ modulo $\mathfrak p$ sits inside
$\mathcal O_L/P$.
\end{hint}
\begin{hint}{ex:qrsplit}
(a) $\operatorname{Gal}(\Q(\zeta_q)/\Q)$ is cyclic of order $q-1$. (b) Theorem~\ref{thm:kd} for $X^2-X-\frac{q^*-1}4$.
(c) $p$ splits in the subfield iff $\operatorname{Frob}_p=\sigma_p$ is trivial on it, iff $\sigma_p$ lies in the unique
subgroup of index $2$.
\end{hint}
\begin{hint}{ex:stickelbergerp}
As in Theorem~\ref{thm:ramdisc}, $d_K\bmod p$ is the discriminant of the trace form of $\mathcal O/p\mathcal O\cong\prod
\mathbb F_{p^{f_i}}$, up to a non-zero square. For $\mathbb F_{p^f}=\mathbb F_p(\beta)$, the discriminant is $\delta^2$
with $\delta=\prod_{i<j}(\beta^{p^i}-\beta^{p^j})$; apply $x\mapsto x^p$ to $\delta$.
\end{hint}

\section{Solutions}

\begin{solution}{ex:sqrtm10}
\idea Kummer--Dedekind with $\theta=\sqrt{-10}$ gives the prime ideals, and a non-principal ideal is detected by the norm
equation $a^2+10b^2=N(P)$.
\writeup Since $-10\equiv2\pmod 4$, $\mathcal O=\Z[\sqrt{-10}]$ (Tome~1, Proposition~11.2), and $N(a+b\sqrt{-10})=a^2+10b^2$.
The four numbers are irreducible: a proper factor of $2$, $7$ or $2\pm\sqrt{-10}$ (of norms $4$, $49$, $14$) would have
norm $2$ or $7$, and $a^2+10b^2\in\{2,7\}$ has no solution.

Modulo $2$, $X^2+10\equiv X^2$: $2\mathcal O=P^2$, $P=(2,\sqrt{-10})$. Modulo $7$, $X^2+10\equiv X^2-4=(X-2)(X+2)$:
$7\mathcal O=QQ'$ with $Q=(7,\sqrt{-10}-2)$ and $Q'=(7,\sqrt{-10}+2)$. The ideal $(2-\sqrt{-10})$ has norm $14$, so it is
the product of one prime of norm $2$ and one of norm $7$. It is contained in $P$ (as $2,\sqrt{-10}\in P$) and in $Q$ (as
$2-\sqrt{-10}=-(\sqrt{-10}-2)$), so it is $PQ$. Likewise $(2+\sqrt{-10})=PQ'$. Then
\[
(2)(7)=P^2\cdot QQ'=(PQ)(PQ')=(2-\sqrt{-10})(2+\sqrt{-10}):
\]
the same four prime ideals, grouped in two ways. If $P=(\alpha)$, then $|N\alpha|=N(P)=2$, impossible; but $P^2=(2)$ is
principal.
\end{solution}

\begin{solution}{ex:cubic}
\idea By Example~\ref{exa:cuberoot2} and Theorem~\ref{thm:kd}, everything is decided by the roots of $X^3-2$ in
$\mathbb F_p$: a root gives a linear factor, and the remaining quadratic factor splits or not.
\writeup If $p\equiv2\pmod3$, $\gcd(3,p-1)=1$, so $x\mapsto x^3$ is injective on the cyclic group $\mathbb F_p^*$, hence
bijective: $2$ has exactly one cube root $c$, and $X^3-2=(X-c)(X^2+cX+c^2)$ where the quadratic factor has no root
(another root would be another cube root of $2$). So $p=P_1P_2$, $f_1=1$, $f_2=2$.

If $p\equiv1\pmod 3$, $\mathbb F_p^*$ contains an element $\omega$ of order $3$. If $c^3=2$, the roots are $c,c\omega,c
\omega^2$, three distinct roots, and $p$ splits completely. If $2$ is not a cube, $X^3-2$ has no root, and a cubic without
roots is irreducible: $p$ is inert. Example~\ref{exa:cubesplit} shows $p=31$ (split) and $p=7$ (inert).
\end{solution}

\begin{solution}{ex:zeta8}
\idea Theorem~\ref{thm:cyclosplit}: in $\Q(\zeta_8)$ only $p\bmod 8$ matters, and the Frobenius $\sigma_p$ is trivial on
the subfield fixed by the subgroup it lies in.
\writeup For odd $p$, $p^2\equiv1\pmod 8$, so $f=1$ if $p\equiv1\pmod 8$ and $f=2$ otherwise; then $g=4/f$. So $p\equiv1
\pmod8$ splits completely into four primes, and the others split into two primes of degree $2$; no odd prime is inert,
as the group $(\Z/8\Z)^*$ is not cyclic. Since $\Phi_8=X^4+1$ and $\mathcal O\supseteq\Z[\zeta_8]$ has index prime to odd
$p$, Theorem~\ref{thm:kd} says $X^4+1$ factors modulo $p$ into $g\ge2$ factors (and modulo $2$ it is $(X+1)^4$).

The subfields: $\sigma_a$ fixes $i=\zeta_8^2$ iff $a\equiv1\pmod4$, fixes $\sqrt2=\zeta_8+\zeta_8^{-1}$ iff $a\equiv\pm1
\pmod8$, and fixes $\sqrt{-2}=\zeta_8+\zeta_8^3$ iff $a\equiv1,3\pmod 8$. Since $p$ splits in a subfield iff $\sigma_p$ is
trivial on it (Exercise~\ref{ex:frobsub}), $p$ splits in $\Q(i)$ iff $p\equiv1\pmod 4$, in $\Q(\sqrt2)$ iff $p\equiv\pm1
\pmod 8$, and in $\Q(\sqrt{-2})$ iff $p\equiv1,3\pmod8$. These are the supplementary laws for $\left(\frac{-1}p\right)$,
$\left(\frac2p\right)$ and $\left(\frac{-2}p\right)$. In each case, $p$ splits in at least one of the three, which is
another way of seeing that $X^4+1$ is never irreducible modulo $p$.
\end{solution}

\begin{solution}{ex:s5}
\idea In characteristic $5$, translation by $1$ permutes the roots, so the Frobenius cannot fix a proper subset unless it
is invariant under translation.
\writeup Let $\beta$ be a root of $f=X^5-X-1$ in a splitting field over $\mathbb F_5$. Since $(\beta+a)^5=\beta^5+a$ for
$a\in\mathbb F_5$, $f(\beta+a)=\beta^5+a-\beta-a-1=0$: the roots are $\beta+a$, $a\in\mathbb F_5$. None is in $\mathbb F_5$,
since $x^5=x$ there, so $f(x)=-1$. So $\beta^5\ne\beta$, and $\beta^5=\beta+c$ for some $c\ne0$ (as $\beta^5$ is a root).
Then $\beta^{5^k}=\beta+kc$, and the first $k\ge1$ with $\beta^{5^k}=\beta$ is $k=5$: $\beta$ has degree $5$ over
$\mathbb F_5$ and $f$ is irreducible.

Modulo $2$: $(X^2+X+1)(X^3+X^2+1)=X^5+2X^4+2X^3+2X^2+X+1\equiv X^5+X+1=f$. Both factors are irreducible (no roots in
$\mathbb F_2$). And $2869=19\cdot151$, with $151$ prime. The same computation as in Exercise~\ref{ex:discbinom} (with
$f'(\theta_i)=5\theta_i^4-1$) gives $\operatorname{disc}(X^5+aX+b)=4^4a^5+5^5b^4$, which is $-256+3125=2869$ here.
\end{solution}

\begin{solution}{ex:twogen}
\idea Only the primes dividing $a$ matter: away from them, $a\mathcal O$ already gives everything. Choose $b$ with the
exact right valuation at each of them.
\writeup Let $P_1,\dots,P_s$ be the primes dividing $a\mathcal O$, and $k_i=v_{P_i}(I)$. Since $a\in I$,
$v_P(I)\le v_P(a)$ for all $P$, so $v_P(I)=0$ for $P\ne P_i$: $I=\prod_iP_i^{k_i}$. For each $i$ pick $b_i\in
P_i^{k_i}\setminus P_i^{k_i+1}$ (these ideals differ by unique factorisation). The ideals $P_i^{k_i+1}$ are pairwise coprime,
so by the Chinese remainder theorem there is $b\in\mathcal O$ with $b\equiv b_i\pmod{P_i^{k_i+1}}$ for every $i$. Then
$b\in P_i^{k_i}$ and $b\notin P_i^{k_i+1}$, that is $v_{P_i}(b)=k_i$; so $b\in I$. By Corollary~\ref{cor:divides}(c),
$a\mathcal O+b\mathcal O$ has exponent $\min(v_P(a),v_P(b))$ at $P$: at $P_i$ it is $\min(v_{P_i}(a),k_i)=k_i$, and at
any other $P$ it is $\min(0,v_P(b))=0$. So $a\mathcal O+b\mathcal O=\prod P_i^{k_i}=I$.
\end{solution}

\begin{solution}{ex:totram}
\idea The units $\frac{1-\zeta^a}{1-\zeta}$ make all the $1-\zeta^a$ ``the same'' up to units, and their product is
$\Phi_q(1)=\ell$.
\writeup The primitive $q$-th roots of unity are the roots of $X^{q}-1$ that are not roots of $X^{q/\ell}-1$, so $\Phi_q=
\frac{X^q-1}{X^{q/\ell}-1}=\Phi_\ell(X^{q/\ell})=1+X^{q/\ell}+\dots+X^{(\ell-1)q/\ell}$, and $\Phi_q(1)=\ell$. So
$\ell=\prod_{\gcd(a,\ell)=1}(1-\zeta^a)$. For $\gcd(a,\ell)=1$, $u_a=\frac{1-\zeta^a}{1-\zeta}=1+\zeta+\dots+\zeta^{a-1}\in
\Z[\zeta]$, and with $b$ such that $ab\equiv1\pmod q$, $u_a^{-1}=\frac{1-(\zeta^a)^b}{1-\zeta^a}\in\Z[\zeta]$ too. So
$u_a$ is a unit and $\ell=(1-\zeta)^{\varphi(q)}\cdot u$ with $u$ a unit: $\ell\mathcal O=(1-\zeta)^{\varphi(q)}\mathcal O$.
By Theorem~\ref{thm:sumef}, $\sum e_if_i=\varphi(q)$; the ideal $(1-\zeta)$ is a product of primes, and its
$\varphi(q)$-th power is $\ell\mathcal O$; comparing exponents, $(1-\zeta)$ must be a single prime with $e=\varphi(q)$ and
$f=1$. (Directly: $N(1-\zeta)=\Phi_q(1)=\ell$ is prime, so $(1-\zeta)$ is a prime of norm $\ell$.)
\end{solution}

\begin{solution}{ex:dedekindfield}
\idea In $\mathbb F_2^3$ every element satisfies $x^2=x$. Conversely, a ring of order $8$ where all elements are
idempotent is $\mathbb F_2^3$. Then count: three primes of degree $1$ above $2$, but $\mathbb F_2[X]$ has only two monic
polynomials of degree $1$.
\writeup From $\beta=\frac{\theta^2+\theta}2$, $\theta^2=2\beta-\theta$. Using $\theta^3=\theta^2+2\theta+8$ one computes
$\theta\beta=\frac{\theta^3+\theta^2}2=\theta^2+\theta+4=4+2\beta$, and $\beta^2=\frac{\theta(\theta\beta)+\theta\beta}{2}
=\frac{(\theta+1)(4+2\beta)}2=(\theta+1)(2+\beta)=2\theta+2+\theta\beta+\beta=6+2\theta+3\beta$. Modulo $2\mathcal O$:
$\theta^2\equiv\theta$, $\beta^2\equiv\beta$, $\theta\beta\equiv0$. For $x=a+b\theta+c\beta$, $x^2\equiv a^2+b^2\theta^2+
c^2\beta^2\equiv a+b\theta+c\beta=x$ (cross terms are even, and $a^2\equiv a$). So $e_1=\bar\theta$, $e_2=\bar\beta$,
$e_3=1-\bar\theta-\bar\beta$ are idempotents with $e_ie_j=0$ for $i\ne j$ (for instance $e_1e_3=\bar\theta-\bar\theta^2-
\bar\theta\bar\beta=0$) and $e_1+e_2+e_3=1$. The map $x\mapsto(xe_1,xe_2,xe_3)$ is then a ring isomorphism
$\mathcal O/2\mathcal O\cong\mathbb F_2e_1\times\mathbb F_2e_2\times\mathbb F_2e_3\cong\mathbb F_2^3$ (each $xe_i$ is $0$ or
$e_i$, as $\mathcal O/2\mathcal O$ is spanned by $1,\bar\theta,\bar\beta$). This ring has three maximal ideals, each with
quotient $\mathbb F_2$, and no nilpotents: $2\mathcal O=P_1P_2P_3$ with $e_i=f_i=1$.

If $\mathcal O=\Z[\gamma]$, Step~1 of Theorem~\ref{thm:kd} (with index $1$) gives $\mathcal O/2\mathcal O\cong\mathbb
F_2[X]/(\bar m_\gamma)$, whose maximal ideals correspond to the distinct irreducible factors of $\bar m_\gamma$, with
residue degree their degree. Three factors of degree $1$ would need three distinct monic polynomials of degree $1$ over
$\mathbb F_2$, and there are only $X$ and $X+1$. So no $\gamma$ generates $\mathcal O$: $2$ is a \emph{common index
divisor} of $K$.
\end{solution}

\begin{solution}{ex:closure}
\idea For the Galois field $L$, splitting completely means ``the Frobenius is trivial''. Proposition~\ref{prop:frobembed}
transports this to $K$ through its embeddings. Ramification must be handled separately, and nilpotents detect it.
\writeup \emph{Ramification.} For any number field $F$, $p$ is unramified in $F$ iff $\mathcal O_F/p\mathcal O_F\cong
\prod\mathcal O_F/P_i^{e_i}$ has no non-zero nilpotent element: if some $e_i\ge2$, the proof of Theorem~\ref{thm:ramdisc}
exhibits one; if all $e_i=1$ the ring is a product of fields. Now if $x\in\mathcal O_K\setminus p\mathcal O_K$ and $x^k\in
p\mathcal O_K$, then $x^k\in p\mathcal O_L$, while $x\notin p\mathcal O_L$ (otherwise $\frac xp\in\mathcal O_L\cap K=
\mathcal O_K$). So if $p$ ramifies in $K$, it ramifies in $L$; in that case it splits completely in neither, and the
claim holds. From now on $p$ is unramified in $K$. Let $P$ be a prime of $L$ above $p$ and $\sigma\in D_P$ lifting the
Frobenius of $\mathcal O_L/P$ (Theorem~\ref{thm:decomp}).

\emph{If $p$ splits completely in $L$}, then $D_P$ has order $ef=1$, so $\sigma=1$, and by
Proposition~\ref{prop:frobembed} all the $f_i$ are $1$: $p$ splits completely in $K$.

\emph{If $p$ splits completely in $K$}, all $f_i=1$, so by Proposition~\ref{prop:frobembed} $\sigma\circ\tau=\tau$ for
every embedding $\tau:K\to L$: $\sigma$ is the identity on every conjugate field $\tau(K)$. These fields generate $L$
(the Galois closure is generated by the conjugates of a primitive element of $K$), so $\sigma=1$. Its reduction, the
Frobenius of $\mathcal O_L/P$, is then trivial, so $f(P)=1$. But then $\mathcal O_L/P=\mathbb F_p$, and \emph{every}
element of $D_P$ reduces to the identity, which is the Frobenius: the argument applies to every $\sigma\in D_P$, and
$D_P=1$. So $e(P)f(P)=|D_P|=1$ and $p$ splits completely in $L$.
\end{solution}

\begin{solution}{ex:frobsub}
\idea Both sides are elements of $\operatorname{Gal}(K/\Q)$ that fix $\mathfrak p$ and reduce to $x\mapsto x^p$; there
is only one such.
\writeup Let $\sigma=\operatorname{Frob}_P$ and $\sigma'=\sigma|_K$, an automorphism of $K$ since $K/\Q$ is Galois. For
$x\in\mathcal O_K$, $\sigma'(x)-x^p=\sigma(x)-x^p\in P\cap\mathcal O_K=\mathfrak p$. If $x\in\mathfrak p$, then
$\sigma'(x)\equiv x^p\equiv0$, so $\sigma'(\mathfrak p)\subseteq\mathfrak p$ and $\sigma'\in D_{\mathfrak p}$, with
reduction $x\mapsto x^p$. Since $p$ is unramified in $L$, it is unramified in $K$ ($e$ is multiplicative in towers),
so $D_{\mathfrak p}\to\operatorname{Gal}(k_{\mathfrak p}/\mathbb F_p)$ is injective (Theorem~\ref{thm:decomp}) and
$\sigma'=\operatorname{Frob}_{\mathfrak p}$.
\end{solution}

\begin{solution}{ex:qrsplit}
\idea Compute the splitting of $p$ in the quadratic field in two ways: by a polynomial modulo $p$ (giving
$\left(\frac{q^*}p\right)$), and by the Frobenius $\sigma_p$ in the cyclotomic field (giving $\left(\frac pq\right)$).
\writeup (a) $G=\operatorname{Gal}(\Q(\zeta_q)/\Q)\cong(\Z/q\Z)^*$ is cyclic of even order $q-1$, so it has exactly one
subgroup $H$ of index $2$, the squares, and by Galois theory exactly one quadratic subfield $F=\Q(\zeta_q)^H$. The Gauss
sum $g=\sum_a\left(\frac aq\right)\zeta_q^a$ satisfies $g^2=q^*$ (Tome~1, Exercise~6.7), so $\Q(\sqrt{q^*})\subseteq
\Q(\zeta_q)$ is quadratic, and it is $F$.

(b) $q^*\equiv1\pmod 4$, so $\mathcal O_F=\Z[\omega]$ with $\omega=\frac{1+\sqrt{q^*}}2$, root of $X^2-X-\frac{q^*-1}4$, of
discriminant $q^*$. For odd $p\ne q$, Theorem~\ref{thm:kd} applies, and the polynomial has a root modulo $p$ iff its
discriminant $q^*$ is a square modulo $p$ ($p$ odd: complete the square). So $p$ splits iff $\left(\frac{q^*}p\right)=1$,
and otherwise is inert.

(c) By Exercise~\ref{ex:frobsub}, $\operatorname{Frob}_p$ in $F$ is the restriction of $\sigma_p$, and $p$ splits in $F$
iff it is trivial there, iff $\sigma_p\in H$, iff $p$ is a square modulo $q$.

(d) So $\left(\frac{q^*}p\right)=\left(\frac pq\right)$. With $\left(\frac{q^*}p\right)=\left(\frac{-1}p\right)^{(q-1)/2}
\left(\frac qp\right)=(-1)^{\frac{p-1}2\frac{q-1}2}\left(\frac qp\right)$, this is quadratic reciprocity.
\end{solution}

\begin{solution}{ex:stickelbergerp}
\idea Reduce modulo $p$ as in Theorem~\ref{thm:ramdisc}: $\left(\frac{d_K}p\right)$ only depends on the discriminant of
the trace form of $\mathcal O/p\mathcal O$, a product of finite fields. For a single finite field $\mathbb F_{p^f}$,
the square root $\delta$ of the discriminant is a Vandermonde product, and the Frobenius permutes its factors by an
$f$-cycle, of sign $(-1)^{f-1}$; so $\delta$ is in $\mathbb F_p$ exactly when $f$ is odd.
\writeup For a finite-dimensional commutative $\mathbb F_p$-algebra $B$ with non-degenerate trace form, let
$\Delta(B)\in\mathbb F_p^*$ be the determinant of $(\operatorname{Tr}_B(b_ib_j))$ in a basis $(b_i)$. A change of basis
multiplies it by a non-zero square, so $\chi(B)=\left(\frac{\Delta(B)}p\right)$ is well defined, and $\chi(B_1\times
B_2)=\chi(B_1)\chi(B_2)$ (block-diagonal matrix in an adapted basis). As in the proof of Theorem~\ref{thm:ramdisc}, $d_K
\bmod p=\Delta(\mathcal O/p\mathcal O)$, and $\mathcal O/p\mathcal O\cong\prod_{i=1}^g\mathbb F_{p^{f_i}}$ because $p$ is
unramified. So $\left(\frac{d_K}p\right)=\prod_i\chi(\mathbb F_{p^{f_i}})$.

Let $F=\mathbb F_{q}$, $q=p^f$, $F=\mathbb F_p(\beta)$, with basis $1,\beta,\dots,\beta^{f-1}$, and $\beta_k=\beta^{p^k}$,
$0\le k<f$, the conjugates of $\beta$. As in the proof of Theorem~\ref{thm:ramdisc}, $\operatorname{Tr}_F(u)=\sum_ku^{p^k}$,
so $\operatorname{Tr}_F(\beta^{i+j})=\sum_k\beta_k^i\beta_k^j$, and the matrix of the trace form is $V^{\mathsf T}V$ with
$V=(\beta_k^{\,j})_{k,j}$ a Vandermonde matrix. Hence $\Delta(F)=\delta^2$ with $\delta=\det V=\pm\prod_{i<j}(\beta_j-
\beta_i)\ne0$. The Frobenius $\varphi(x)=x^p$ sends $\beta_k$ to $\beta_{k+1}$ (indices modulo $f$), so applying it to
the entries of $V$ permutes its rows by an $f$-cycle, of sign $(-1)^{f-1}$; so $\varphi(\delta)=(-1)^{f-1}\delta$. If $f$
is odd, $\delta$ is fixed by $\varphi$, so $\delta\in\mathbb F_p$ and $\Delta(F)=\delta^2$ is a square: $\chi(F)=1$. If
$f$ is even, $\varphi(\delta)=-\delta\ne\delta$ ($p$ is odd), so $\delta\notin\mathbb F_p$; and $\Delta(F)$ is not a square
in $\mathbb F_p$, since its only square roots in $F$ are $\pm\delta$. So $\chi(F)=-1$. In both cases
$\chi(\mathbb F_{p^f})=(-1)^{f-1}$, and
\[
\left(\frac{d_K}p\right)=\prod_{i=1}^g(-1)^{f_i-1}=(-1)^{\sum f_i-g}=(-1)^{n-g}.
\]
Check on $\Q(\sqrt[3]2)$, $d_K=-108$, with Example~\ref{exa:cubesplit}: for $p=5$, $g=2$ and $\left(\frac{-108}5\right)=
\left(\frac25\right)=-1=(-1)^1$. For $p=7$, $g=1$, and $-108\equiv4$ is a square modulo $7$. For $p=31$, $g=3$, and
$-108\equiv16$ is a square modulo $31$.
\end{solution}

\section*{Sources and further reading}

\begin{itemize}
\item D.~A.~Marcus, \emph{Number Fields}, 2nd ed., Universitext, Springer, 2018, Chapters~3--4. Unique factorisation,
Kummer--Dedekind, and the Galois theory of primes, with a wealth of exercises (Exercise~\ref{ex:qrsplit} is in Chapter~4).
\item J.~Neukirch, \emph{Algebraic Number Theory}, Grundlehren~322, Springer, 1999, Chapter~I, \S\S3, 8--9. The same
theory for Dedekind domains, with the geometric picture of Figure~\ref{fig:spec}.
\item D.~Mumford, \emph{The Red Book of Varieties and Schemes}, Lecture Notes in Math.~1358, Springer, 1988, Chapter~II,
\S1. The original drawing of $\operatorname{Spec}\Z[X]$ that inspired Figure~\ref{fig:spec}.
\item E.~Noether, ``Abstrakter Aufbau der Idealtheorie in algebraischen Zahl- und Funktionenk\"orpern'',
\emph{Math.\ Ann.}\ 96 (1927), 26--61. The axioms behind Proposition~\ref{prop:noeth} and Theorem~\ref{thm:ufi}.
\item R.~Dedekind, ``Konstruktion von Quaternionk\"orpern'' (1878--1882, published posthumously in his \emph{Gesammelte
Werke}, vol.~2), where Theorem~\ref{thm:dedekindcycle} appears; see also Tome~2, Chapter~10.
\item P.~Stevenhagen and H.~W.~Lenstra, Jr., ``Chebotar\"ev and his density theorem'', \emph{Math.\ Intelligencer}
18 (1996), 26--37. The history behind Figure~\ref{fig:cycletypes}; we follow it in Chapter~8.
\item H.~Cohen, \emph{A Course in Computational Algebraic Number Theory}, GTM~138, Springer, 1993, \S\S3.4 and~6.2.
Factoring polynomials modulo $p$, and decomposing primes in practice.
\item \emph{History.} The sessions of the Paris Academy in March--May 1847 (Lam\'e, Liouville, Cauchy's sealed envelopes,
Kummer's letter read by Liouville on 24 May) are told from the \emph{Comptes rendus} in H.~M.~Edwards, \emph{Fermat's Last
Theorem}, Springer GTM~50, 1977, Chapter~4. Kummer's 1844 failure of unique factorisation for $\lambda=23$ appeared in his
Breslau memoir of that year, reprinted in \emph{J.\ Math.\ Pures Appl.}\ 12 (1847), 185--212.
\end{itemize}
